\section{测度把长度、计数与概率统一起来}
\label{sec:12-Real-Analysis-Support/05-Measure-Integration-and-Conditional-Expectation/02}

一份故障复盘里并排放着三个数：告警持续了 47 分钟，受影响的用户有 12304 个，随机一次请求命中故障的概率是 \(3\%\)。单位互不相同，一个是时间，一个是人头，一个没有量纲。可是三个数遵守同一条算术规则：把对象拆成互不重叠的几块，各块的数加起来，等于整体的数。

时长、计数、概率，三套语言各自有一套定理，读起来像三门课。它们服从同一条加法律这件事，值得单独写下来当作一个对象——一旦写下来，为长度证明的极限定理不必为概率再证一遍。

\subsection{三种大小服从同一条加法律}

\begin{definition}
设 \((\Omega,\mathcal F)\) 是可测空间。函数 \(\mu:\mathcal F\to[0,\infty]\) 称为\textbf{测度}，如果 \(\mu(\varnothing)=0\)，且对任意两两不交的 \(A_1,A_2,\ldots\in\mathcal F\) 有
\[
\mu\Bigl(\bigcup_{n=1}^{\infty}A_n\Bigr)=\sum_{n=1}^{\infty}\mu(A_n).
\]
三元组 \((\Omega,\mathcal F,\mu)\) 称为测度空间。
\end{definition}

允许取值 \(\infty\) 不是技术上的将就。整条实数轴的长度、可数无限集合的元素个数，本来就该是无穷；强行排除它们，反而要在每条定理里补一堆有界性假设。

复盘里那三个数分别对应三个测度。计数测度令 \(\mu(A)\) 等于 \(A\) 的元素个数，无限集取 \(\infty\)；沿它做积分就退化成求和。Lebesgue 测度 \(\lambda\) 在实轴上由 \(\lambda((a,b))=b-a\) 出发，延拓到全体 Borel 集。概率测度是额外要求 \(P(\Omega)=1\) 的测度，归一化不改动任何结构，只把总质量钉在 1 上。

这条统一带来的实惠很直接：关于积分与极限交换的所有结论，只要在测度层面证过一次，就同时是关于求和的、关于长度的、关于期望的结论。离散与连续之间那道教科书式的分界线，在这一层已经不存在了。

\subsection{有限可加撑不到极限}

定义里写的是可数可加。为什么不满足于有限可加——不交的两块能相加，三块也能，逐次推下去不就够了吗。

\begin{center}
\begin{tikzpicture}[>=stealth]
\draw[thick] (0,1.6) -- (8,1.6);
\foreach \x in {0,4,6,7,7.5,7.75,8} \draw[thick] (\x,1.45) -- (\x,1.75);
\node at (2,2.0) {\footnotesize \(I_1\)};
\node at (5,2.0) {\footnotesize \(I_2\)};
\node at (6.5,2.0) {\footnotesize \(I_3\)};
\node at (7.4,2.0) {\footnotesize \(\cdots\)};
\node at (0,1.05) {\footnotesize \(0\)};
\node at (8,1.05) {\footnotesize \(1\)};
\node at (4,2.55) {\footnotesize 可数多段互不重叠，长度和恰为 \(1\)};
\draw[thick] (0,0.1) -- (7,0.1);
\draw[thick,gray] (7,0.1) -- (8,0.1);
\foreach \x in {0,4,6,7} \draw[thick] (\x,-0.05) -- (\x,0.25);
\draw[gray] (8,-0.05) -- (8,0.25);
\draw[<->,gray] (7,-0.45) -- (8,-0.45);
\node[gray] at (7.5,-0.85) {\footnotesize 缺口 \(1/8\)};
\node at (3,-0.45) {\footnotesize 只取前三段：永远差一截};
\end{tikzpicture}
\end{center}

\emph{有限可加只管得住任意有限步，可数可加才把极限位置上的那一截缺口一并算清。}

把 \([0,1]\) 劈成 \(I_1=[0,1/2)\)、\(I_2=[1/2,3/4)\)、\(I_3=[3/4,7/8)\)，如此以下。有限可加性告诉你前 \(n\) 段的长度和是 \(1-2^{-n}\)，对每个 \(n\) 都对，但它从不允许你把这句话在 \(n\to\infty\) 处收口。要断言这可数多段确实穷尽了整个区间、长度和正好是 1，得先有可数可加性。

概率论里这一步不是修饰。\hyperref[sec:05-Probability/01-Probability-Spaces/02]{``概率公理与概率空间''}把可数可加直接写进公理，因为但凡涉及极限的陈述——大数定律里的几乎必然收敛、``某事件发生无穷多次''、首次击中时间的分布——都要求在可数并处成立等式，而不只是在每个有限截断处成立。用有限可加性构建的概率论，能算清任何一次有限试验，却说不出任何一句关于长期行为的话。

\subsection{连续性是同一条性质的另一副面孔}

可数可加性可以改写成关于集合列极限的连续性，实际使用中多半以这副面孔出现。

设 \(A_1\subseteq A_2\subseteq\cdots\)，则
\[
\mu\Bigl(\bigcup_{n=1}^{\infty}A_n\Bigr)=\lim_{n\to\infty}\mu(A_n).
\]
骨架是把递增列拆成不交增量：令 \(B_1=A_1\)、\(B_n=A_n\setminus A_{n-1}\)，则诸 \(B_n\) 互不相交，前 \(n\) 个的并等于 \(A_n\)，对 \(\bigcup_n B_n\) 用可数可加即得。含义是集合缓缓涨上去时，质量不会在极限那一刻凭空冒出来。

反方向要多一个条件。设 \(A_1\supseteq A_2\supseteq\cdots\) 且 \(\mu(A_1)<\infty\)，则
\[
\mu\Bigl(\bigcap_{n=1}^{\infty}A_n\Bigr)=\lim_{n\to\infty}\mu(A_n).
\]
有限性删不得。取 \(A_n=(n,\infty)\)，每个 \(A_n\) 的 Lebesgue 测度都是 \(\infty\)，而 \(\bigcap_n A_n=\varnothing\) 测度为零。补集论证走到``\(\infty-\infty\)''就断了，那个式子提供不了任何控制。概率测度总质量为 1，这个条件自动满足，所以概率论里从上连续用得毫无顾忌；一旦换到无穷测度空间，就得回头检查。

\subsection{不可测集是这条加法律的定价}

上一节提过 Vitali 的构造。现在可以说清它到底在为什么付账。

骨架是这样：在 \([0,1)\) 上把相差有理数的点归为一类，用选择公理从每一类里挑一个代表，凑成集合 \(V\)。把 \(V\) 沿全体有理数 \(q\in\Q\cap[-1,1)\) 平移，得到可数多个互不相交的副本，它们的并夹在 \([0,1)\) 与 \([-1,2)\) 之间。平移不变性要求每个副本长度相同，记作 \(c\)。若 \(c=0\)，可数可加给出总长度 0，与下界 1 矛盾；若 \(c>0\)，可数多个相同正数相加得 \(\infty\)，与上界 3 矛盾。两条路都堵死，结论是 \(V\) 压根不能被赋予长度。

矛盾出在哪里，值得看仔细：出在``可数多个相同的数相加''这一步。若只要求有限可加，这个论证走不下去——Banach 在 1923 年证明，直线与平面上确实存在定义在全部子集上、平移不变的有限可加集函数。换句话说，不可测集不是实数轴的什么病灶，是可数可加性开出的账单。想要极限定理，就得放弃``每个集合都有长度''。三维以上连有限可加的版本也保不住，Banach--Tarski 分球说的就是这件事，代价随维数还在涨。

选择公理在这条链子上是不可省的一环：不用它，无法构造 \(V\)。存在一些集合论模型，其中所有实数子集都 Lebesgue 可测。现代概率论选择了留下选择公理、承认不可测集存在，然后从一开始就把讨论限定在 \(\sigma\)-代数内。上一节强调``先定 \(\mathcal F\) 再定 \(\mu\)''，理由到这里才完整。

\subsection{零测集撑起几乎处处这个说法}

单点的 Lebesgue 测度为零。有理数集是可数多个单点的并，可数可加立刻给出
\[
\lambda(\Q)=0,
\]
尽管 \(\Q\) 在每个开区间里都稠密。反过来，不可数的零测集也存在，Cantor 集就是。基数与测度回答的是两个问题：一个问点有多少，另一个问按指定的大小规则占了多少质量，两者不互相决定。

性质若在某个零测集之外处处成立，就说它几乎处处成立。两个函数可以在一个稠密的零测集上取值不同，而作为积分对象完全等价——这正是 \(L^p\) 空间的元素是等价类而不是单个函数的原因，见 \hyperref[sec:12-Real-Analysis-Support/07-Lp-Spaces-and-Function-Geometry/01]{``\(L^p\) 范数把积分变成距离''}。

零测集还牵出一个技术细节。若 \(N\in\mathcal F\)、\(\mu(N)=0\)，而 \(N\) 的某个子集不在 \(\mathcal F\) 中，这个测度空间叫作不完备的。把所有零测集的子集补进事件族并令其测度为零，得到完成化。Borel \(\sigma\)-代数配 Lebesgue 测度尚未完成，完成之后才是 Lebesgue 可测集族。概率论通常默认工作在完成化后的空间上，但涉及可测性的具体论证仍要分清是 Borel 可测还是 Lebesgue 可测，两个族并不一样大。

\subsection{分布首先是测度，密度是后话}

可测映射能把测度搬家。若 \(X:(\Omega,\mathcal F)\to(S,\mathcal S)\) 可测，定义
\[
\mu_X(B)=\mu\bigl(X^{-1}(B)\bigr),\qquad B\in\mathcal S,
\]
称为 \(X\) 的推前测度。可测性保证右边有意义，可数可加性顺着原像运算原样传过去。概率情形下 \(\mu_X\) 就是 \(X\) 的分布，写作 \(P_X(B)=P(X\in B)\)。

这条定义里没有出现概率质量函数，也没有出现密度。离散分布、连续分布、离散连续混合的分布、集中在低维曲面上的分布，一律先是测度。``它有没有密度''是一个额外的问题，而且只有指明相对于哪个参考测度才提得出来——第四节专门处理它。眼下更要紧的是：手里有了一个可数可加的非负集函数，怎样沿着它做积分。
